\section*{Problem 2}
% \paragraph{Problem.}
In an abstract sequence of length $18$, we are given candidate exons (closed intervals) with weights:
\[
    [1,4]:4,\quad [3,6]:5,\quad [3,8]:10,\quad [9,10]:1,\quad [8,16]:12,\quad [15,18]:6,\quad [13,16]:7,\quad [7,11]:8.
\]
Find a maximum-weight subset of \emph{non-overlapping} exons. We treat intervals as overlapping if they share at least one position; thus two chosen intervals $[L_i,R_i]$ and $[L_j,R_j]$ must satisfy $R_i<L_j$ or $R_j<L_i$.

\paragraph{Model.}
This is the classic \emph{Weighted Interval Scheduling} (WIS) problem. Let each exon $j$ have left endpoint $L_j$, right endpoint $R_j$, and weight $w_j$.
Sort intervals by nondecreasing $R$. For each interval $j$, define the \emph{compatibility index}
\[
    p(j)\;=\;\max\{\,i<j:\ R_i<L_j\,\},
\]
with $p(j)=0$ if no such $i$ exists. The optimal value $M[j]$ for the first $j$ intervals satisfies the recurrence
\[
    M[0]=0,\qquad
    M[j]=\max\!\big(M[j-1],\,w_j+M[p(j)]\big)\quad (j\ge 1).
\]
We reconstruct an optimal set by taking interval $j$ exactly when $w_j+M[p(j)]>M[j-1]$.

\paragraph{Data preparation.}
Label the intervals for convenience:
\[
    \begin{array}{lcl}
        A=[1,4]:4,   & \quad & B=[3,6]:5,\quad C=[3,8]:10,\quad D=[9,10]:1,    \\
        E=[8,16]:12, & \quad & F=[15,18]:6,\quad G=[13,16]:7,\quad H=[7,11]:8.
    \end{array}
\]
Sort by $R$:
\[
    A(1,4,4),\; B(3,6,5),\; C(3,8,10),\; D(9,10,1),\; H(7,11,8),\; E(8,16,12),\; G(13,16,7),\; F(15,18,6).
\]
Compute $p(j)$ using the strict non-overlap condition $R_i<L_j$:
\[
    p(1)=0,\; p(2)=0,\; p(3)=0,\; p(4)=3,\; p(5)=2,\; p(6)=2,\; p(7)=5,\; p(8)=5.
\]

\paragraph{Dynamic programming table.}
Let $w_j$ denote the weight of the $j$-th interval in the sorted list. The DP values are:

\medskip
\begin{tabular}{@{}cccccc@{}}
    \toprule
    $j$ & Interval $(L,R,w)$ & $p(j)$ & $w_j+M[p(j)]$ & $M[j-1]$ & $M[j]=\max$ \\
    \midrule
    0   & --                 & --     & --            & --       & $0$         \\
    1   & $(1,4,4)$          & $0$    & $4$           & $0$      & $4$         \\
    2   & $(3,6,5)$          & $0$    & $5$           & $4$      & $5$         \\
    3   & $(3,8,10)$         & $0$    & $10$          & $5$      & $10$        \\
    4   & $(9,10,1)$         & $3$    & $1+10=11$     & $10$     & $11$        \\
    5   & $(7,11,8)$         & $2$    & $8+5=13$      & $11$     & $13$        \\
    6   & $(8,16,12)$        & $2$    & $12+5=17$     & $13$     & $17$        \\
    7   & $(13,16,7)$        & $5$    & $7+13=20$     & $17$     & $20$        \\
    8   & $(15,18,6)$        & $5$    & $6+13=19$     & $20$     & $20$        \\
    \bottomrule
\end{tabular}

\medskip
Thus the optimal value is $M[8]=\boxed{20}$.

% \paragraph{Reconstruction (one optimal solution).}
% From $j=8$: $M[8]=M[7]$ (do not take $F$).
% At $j=7$: $7+M[5]=20>M[6]=17$ (take $G=[13,16]:7$), then jump to $j=p(7)=5$.
% At $j=5$: $8+M[2]=13>M[4]=11$ (take $H=[7,11]:8$), then jump to $j=p(5)=2$.
% At $j=2$: $5+M[0]=5>M[1]=4$ (take $B=[3,6]:5$), then jump to $j=p(2)=0$ and stop.

\paragraph{Answer.}
A maximum-weight non-overlapping set is
\[
    \boxed{\{[3,6],[7,11],[13,16]\}}
\]
with total weight
\[
    \boxed{5+8+7=20}.
\]

% \paragraph{Remarks.}
% A brute-force check would consider all $2^8$ subsets and keep the best feasible one; the DP above runs in $O(n\log n)$ time (for sorting) plus $O(n\log n)$ or $O(n)$ to compute $p(j)$ (by binary search or two-pointer sweep), and $O(n)$ for the DP, which is far more efficient and certifies optimality via the recurrence.